\section{问题三：月度招工与人员分配模型}

\subsection{招工人数下界}

设月度招工人数为$N$。一方面，任意一天到岗人数不能超过总人数，因此$N\ge\max_dr_d=581$；另一方面，每人工作23天，全月可提供$23N$人日，故
\begin{equation}
N\ge\left\lceil\frac{\sum_{d=1}^{30}r_d}{23}\right\rceil
=\left\lceil\frac{11951}{23}\right\rceil=520.
\end{equation}
因此理论下界为
\begin{equation}N_{\mathrm{LB}}=\max\{581,520\}=581.\label{eq:lb}\end{equation}

\subsection{工人--日期排班模型}

为证明下界可达，定义$w_{id}=1$表示工人$i$在第$d$天工作。每人恰好工作23天：
\begin{equation}\sum_{d=1}^{30}w_{id}=23,\qquad i=1,\ldots,N.\end{equation}
“连续工作不超过7天”等价于任意连续8天窗口至少休息1天：
\begin{equation}\sum_{d=t}^{t+7}w_{id}\le7,
\quad t=1,\ldots,23.\label{eq:window}\end{equation}
每日覆盖需求为
\begin{equation}\sum_{i=1}^{N}w_{id}\ge r_d,\qquad d=1,\ldots,30.\label{eq:cover}\end{equation}

从$N=581$开始求解式\eqref{eq:window}--\eqref{eq:cover}，模型含17430个0--1出勤变量。求解得到可行矩阵：每人工作23天，任意连续8天最多工作7天，每日到岗人数均不低于$r_d$。随后将每日到岗工人依次分配给问题二的5个班次；超出最低需求的人员加入当日人数最多的班次作为机动人员。由此构造出完整“工人--日期--班次”方案，证明式\eqref{eq:lb}下界可达，即最少招工581人。

\begin{table}[htbp]\centering
\caption{问题三月度排班可行性检验}\label{tab:roster-check}
\begin{tabular}{lrr}\toprule 指标 & 要求 & 计算结果\\\midrule
招工人数 & 最少 & 581\\
每人工日数 & 23 & 23--23\\
连续工作上限 & 不超过7天 & 7\\
全月总工日 & 不低于需求 & 13363（需求11951）\\
\bottomrule\end{tabular}\end{table}

图\ref{fig:roster}显示每日实际到岗人数与最低需求。由于每名工人必须工作23天，方案总计提供$581\times23=13363$人日，比最低需求多1412人日；这些冗余是固定工日制度与峰值用工共同作用的结果，而非模型浪费。

\begin{figure}[htbp]
  \centering
  \includegraphics[width=.90\textwidth]{figures/monthly_roster.pdf}
  \caption{问题三每日最低需求与实际到岗人数}
  \label{fig:roster}
\end{figure}
